% !TeX program = pdflatex
% =============================================================================
% Arbeitsblatt 3: Kraftaddition & Kraftzerlegung
% UZH Praktikum I -- Sara Romer Urban, Kantonsschule im Lee, HS2026
% Klasse: 2e (02.09.2026)
%
% Direkte Fortsetzung von Gewichtskraft & Federkraft (26.08.2026): dieselbe
% Feder-Masse-Skizze dient hier als Repetitions-/Anknuepfungs-Einstieg, an dem
% zum ersten Mal explizit das Kraeftegleichgewicht als Vektorgleichung
% formuliert wird. Enthaelt alle Papier-Konstruktionsaufgaben der Lektion
% (Repetition, Kraftaddition zweier und mehrerer Kraefte, Kernaufgabe zur
% Winkelabhaengigkeit der Resultierenden, Kraftzerlegung mit dem
% Kraefteparallelogramm) -- ausschliesslich grafisch, keine rechnerische/
% komponentenweise Kraftaddition (siehe Abgrenzung in lernziele.md). Das
% separate Aufgabenblatt mit reinen Uebungsaufgaben entsteht laut
% Lektionsplan Abschnitt 5 als eigenes Dokument.
%
% Quellen: lernziele.md (dieser Ordner), kraefte-addieren-zerlegen-
% lektionsplan.tex (freigegebene Diskussionsgrundlage, dieser Ordner),
% arbeitsblatt2-gewichtskraft-federkraft-teil1.tex (Feder-Masse-Skizze,
% Vorlage fuer Box-Definitionen ohne Farbrahmen), PhysikLibre Kap. 4.1.2
% (Kraft als Vektorgroesse) und 4.3 (4.3.1 Zusammenfassen von Kraeften,
% 4.3.2 Kraefteparallelogramm, 4.3.3 Kraefteediagramm/Kraft-Massstab, 4.3.4
% Aufteilung von Kraeften, Bild 4.17 Ampelbeispiel).
% =============================================================================
\documentclass[addpoints,11pt,a4paper]{exam}

\usepackage[T1]{fontenc}
\usepackage[utf8]{inputenc}
\usepackage[ngerman]{babel}
\usepackage{lmodern}
\usepackage[a4paper,margin=2.2cm,headheight=14pt]{geometry}
\usepackage{amsmath}
\usepackage{siunitx}
% Hausstil: zusammengesetzte Einheiten immer als echter Bruch (\frac), nicht
% als Inline-Schraegstrich oder \tfrac -- gilt global fuer jedes \si/\SI.
\sisetup{per-mode=fraction,fraction-command=\frac}
\usepackage{array,booktabs,tabularx,multirow}
\usepackage{enumitem}
\usepackage{xcolor}
\usepackage{colortbl}
\usepackage{tikz}
\usepackage{physics}
\usepackage[most]{tcolorbox}
\usepackage{lastpage}
\usepackage{graphicx}
\usepackage{hyperref}

\usetikzlibrary{arrows.meta,calc,angles,quotes,decorations.pathmorphing}

\usepackage{setspace}
\usepackage{needspace}
\setlength{\parindent}{0pt}
\setlength{\parskip}{0.6em}
\setstretch{1.08}
\setlist[itemize]{itemsep=0.4em}

% -----------------------------
% Farben (Corporate-Farbschema Physik KFR)
% -----------------------------
\definecolor{titleblue}{HTML}{183A6B}
\definecolor{accentblue}{HTML}{2E6F9E}
\definecolor{lightblue}{HTML}{EAF4FB}
\definecolor{lightgray}{HTML}{F4F6F8}
\definecolor{darkgray}{HTML}{4A4A4A}
\definecolor{accentorange}{HTML}{D9720C}
\definecolor{accentpurple}{HTML}{7B2D8E}
% Hinweisbox: Orange (accentorange), fuer wichtige Klarstellungen.
\definecolor{lightorange}{HTML}{FCEEE0}
% Formelbox: Violett (accentpurple), fuer die zentralen Merksaetze.
\definecolor{lightpurple}{HTML}{F3EAF7}

% Links (hyperref): farblich ans Corporate-Farbschema angeglichen.
\hypersetup{
  colorlinks=true,
  urlcolor=accentblue,
  linkcolor=accentblue,
  pdftitle={Arbeitsblatt 3: Kraftaddition \& Kraftzerlegung},
  pdfauthor={Loris Keller}
}

% -----------------------------
% Spaltentypen
% -----------------------------
\newcolumntype{Y}{>{\centering\arraybackslash}X}
\newcolumntype{P}[1]{>{\raggedright\arraybackslash}p{#1}}

% -----------------------------
% Boxen
% -----------------------------
\tcbset{before skip=10pt, after skip=10pt}
\newtcolorbox{titlebox}{
  enhanced,
  colback=titleblue,
  colframe=titleblue,
  boxrule=0pt,
  arc=3mm,
  left=3mm,right=3mm,top=3mm,bottom=3mm
}

\newlength{\aufgabenindent}
\settowidth{\aufgabenindent}{10.\hskip\labelsep}

% Praktikum I: Lernziele/Einleitung/Theorie/Aufgaben werden NICHT mehr
% farbig gerahmt (nur Formel-, Hinweis- und Antwortbox bleiben Boxen) --
% siehe institutionen/UZH-Praktikum-I/CLAUDE.md, Abschnitt "Boxen-Layout".
\newenvironment{infobox}{%
  \par\addvspace{10pt}\noindent
}{%
  \par\addvspace{10pt}
}

\newenvironment{theoriebox}[1][Theorie]{%
  \par\addvspace{10pt}\noindent\textbf{#1}\par\smallskip\noindent
}{%
  \par\addvspace{10pt}
}

% taskbox/groupbox stehen in questions VOR dem ersten \question (Bilder,
% Fliesstext) -- exam.cls' questions-Liste braucht dafuer eine echte,
% eigenstaendige Box (sonst "Something's wrong--perhaps a missing \item",
% weil z.B. ein \begin{center} vor dem ersten \item der Aussenliste steht).
% Deshalb bleiben es tcolorboxen, nur unsichtbar gemacht (keine Farbe/Rahmen,
% kein Padding) statt sie durch reine Absaetze zu ersetzen.
\newtcolorbox{taskbox}[1][]{
  enhanced, breakable,
  colback=white, colframe=white, boxrule=0pt,
  left=0mm,right=0mm,top=0mm,bottom=0mm,
  before skip=10pt, after skip=10pt,
  before upper={%
    \def\tmpboxtitle{#1}%
    \ifx\tmpboxtitle\empty\else\noindent\textbf{#1}\par\smallskip\fi
  }
}

\newtcolorbox{groupbox}[1][Gruppenarbeit]{
  enhanced, breakable,
  colback=white, colframe=white, boxrule=0pt,
  left=0mm,right=0mm,top=0mm,bottom=0mm,
  before skip=10pt, after skip=10pt,
  before upper={\noindent\textbf{#1}\par\smallskip}
}

% beispielbox: fuer die Einleitung -- wie die anderen Inhaltsboxen ungerahmt.
\newenvironment{beispielbox}[1][Beispiel]{%
  \par\addvspace{10pt}\noindent\textbf{#1}\par\smallskip\noindent
}{%
  \par\addvspace{10pt}
}

% hinweisbox: Orange, fuer wichtige Klarstellungen/Verwechslungswarnungen.
\newtcolorbox{hinweisbox}[1][Hinweis]{
  enhanced,
  breakable,
  colback=lightorange,
  colframe=accentorange,
  title={#1},
  fonttitle=\bfseries,
  boxrule=0.7pt,
  arc=2mm,
  left=5mm,right=5mm,top=2mm,bottom=2mm
}

% formelbox: Violett, um die zentralen Merksaetze dieser Einheit hervorzuheben.
\newtcolorbox{formelbox}[1][Merksatz]{
  enhanced,
  breakable,
  colback=lightpurple,
  colframe=accentpurple,
  title={#1},
  fonttitle=\bfseries,
  boxrule=0.9pt,
  arc=2mm,
  left=5mm,right=5mm,top=2mm,bottom=2mm
}

\newtcolorbox{answerbox}{
  breakable,
  colback=white,
  colframe=gray!55,
  boxrule=0.5pt,
  arc=1.5mm,
  left=2mm,right=2mm,top=2mm,bottom=2mm
}

% Marke: kleines farbiges Inline-Label fuer Teilschritte/Ergebnis-Callouts.
\newtcbox{\Marke}[1][accentpurple]{
  nobeforeafter,
  colback=#1,
  colframe=#1,
  coltext=white,
  boxrule=0pt,
  arc=2pt,
  left=5pt,right=5pt,top=2pt,bottom=2pt,
  fontupper=\bfseries\sffamily\small
}

% --- Feinraster-Helfer fuer alle grafischen Konstruktionsaufgaben ----------
% Feld mit feinem Punkt-/Liniengitter (0.5 cm Nebenlinien, 1 cm Hauptlinien),
% sodass bei Massstab 1 cm = 1 N jede Kraft exakt und nachmessbar
% eingezeichnet werden kann. #1#2 = linke untere Ecke, #3#4 = rechte obere Ecke.
\newcommand{\Feinraster}[4]{%
  \draw[gray!22,step=0.5] (#1,#2) grid (#3,#4);
  \draw[gray!55,step=1] (#1,#2) grid (#3,#4);
}

% Kraft-Massstab-Angabe (z.B. 1 cm = 1 N), einheitlich ueber siunitx statt
% handgesetztem \text{cm}/\text{N} -- #1 = Laenge in cm, #2 = Kraft in N.
\newcommand{\Massstab}[2]{\SI{#1}{\centi\meter}\mathrel{\hat=}\SI{#2}{\newton}}

% --- Loesungen ein-/ausblenden -----------------------------------------------
\noprintanswers
%\printanswers

\pointformat{}
\renewcommand{\solutiontitle}{\noindent\color{titleblue}\textbf{L\"osung:}\enspace}

\pagestyle{headandfoot}
\cfoot{\textcolor{darkgray}{Seite \thepage\ von \pageref{LastPage}}}

\newcommand{\Kopfzeile}[4]{%
  \begin{titlebox}
    {\color{white}\sffamily\bfseries\Large #4}
  \end{titlebox}
  \vspace{0.5em}
  \noindent\textbf{Klasse:}~#1 \hfill \textbf{Datum:}~#2 \hfill \textbf{Thema:}~#3
    \hfill \textbf{Lehrperson:}~Loris Keller
  \vspace{0.8em}
  \lhead{\textcolor{darkgray}{#3}}
  \rhead{\textcolor{darkgray}{#4}}
}

\newlength{\antwortraumlen}
\newcommand{\Antwortraum}[1]{%
  \pgfmathsetlength{\antwortraumlen}{#1*2.5cm}%
  \begin{answerbox}
    \rule{0pt}{\antwortraumlen}%
  \end{answerbox}%
}

% Werte fuer Kopfzeile zentral setzen
\newcommand{\Klasse}{2e}
\newcommand{\Datum}{02.09.2026}
\newcommand{\Thema}{Kraftaddition \& Kraftzerlegung}
\newcommand{\ArbeitsblattTitel}{Arbeitsblatt 3: Kraftaddition \& Kraftzerlegung}

\begin{document}

\Kopfzeile{\Klasse}{\Datum}{\Thema}{\ArbeitsblattTitel}

% =============================================================================
% Lernziele (aus lernziele.md)
% =============================================================================
\begin{infobox}
\textbf{Lernziele}\\
Am Ende dieser Doppellektion kann ich \dots
\begin{itemize}[leftmargin=1.2cm,itemsep=0.3em]
  \item Kraft als vektorielle Grösse (Betrag \emph{und} Richtung) beschreiben und Kräfte an einem Körper korrekt einzeichnen (vom Schwerpunkt aus, mit einheitlichem Kraftmassstab) sowie am Beispiel der Feder-Masse-Skizze das Kräftegleichgewicht als Vektorgleichung formulieren. (LZ1)
  \item die grafische Summe (Resultierende) von zwei und von mehreren an einem Angriffspunkt wirkenden Kräften konstruieren (Kräfteparallelogramm bzw. Kopf-an-Schwanz-Verfahren), mit einheitlichem Kraftmassstab. (LZ2)
  \item an einem selbst konstruierten Beispiel untersuchen, wie der Betrag der resultierenden Kraft von der Winkeldifferenz zweier gleich grosser Kräfte abhängt. (LZ3)
  \item eine gegebene Kraft mithilfe des Kräfteparallelogramms grafisch in zwei vorgegebene Richtungen zerlegen. (LZ4)
  \item für Kraftaddition und Kraftzerlegung je ein Alltagsbeispiel nennen und die grafische Konstruktion darauf anwenden. (LZ5)
\end{itemize}
\end{infobox}

% =============================================================================
% Einleitung
% =============================================================================
\begin{beispielbox}[Einleitung: Wenn mehrere Kräfte gleichzeitig ziehen]
Beim Tauziehen ziehen zwei Mannschaften mit jeweils grosser Kraft am selben
Seil. Trotzdem bewegt sich das Seil oft nur langsam oder gar nicht. Ähnlich
ergeht es einem Boot, das von seinem Motor \emph{und} gleichzeitig vom Wind
seitlich angetrieben wird: Es fährt am Ende weder in die Richtung des Motors
noch in die des Windes, sondern irgendwo dazwischen. In beiden Fällen wirken
mehrere Kräfte gleichzeitig auf denselben Körper, und die Wirkung, die
tatsächlich zu beobachten ist, entspricht keiner der einzelnen Kräfte allein.

Damit man solche Situationen vorhersagen kann, braucht es eine Regel, wie
sich mehrere Kräfte zu einer einzigen \emph{Gesamtwirkung} zusammensetzen,
und ebenso die Umkehrung: Wie ersetzt man eine einzelne Kraft (zum Beispiel
das Gewicht einer an zwei Seilen hängenden Ampel) durch zwei Teilkräfte, die
zusammen dieselbe Wirkung haben? Genau darum geht es in dieser Doppellektion:
Kraftaddition und Kraftzerlegung, beide rein grafisch mit Zeichnung und
Lineal.
\end{beispielbox}

% =============================================================================
% Repetition/Anknuepfung: Feder-Masse-Skizze (LZ1)
% =============================================================================
\begin{questions}
\begin{taskbox}[Repetition: Die Feder-Masse-Skizze]
\question In der letzten Doppellektion haben Sie bereits gezeichnet, dass
sich Federkraft und Gewichtskraft bei einer ruhenden, hängenden Masse die
Waage halten. Die Skizze unten zeigt dieselbe Situation noch einmal.

\begin{center}
\begin{tikzpicture}[>=Latex,scale=1.3]
  \draw[very thick] (-1,3) -- (2,3);
  \foreach \x in {-0.8,-0.4,...,1.8} { \draw (\x,3) -- (\x-0.25,3.3); }
  \draw[decorate,decoration={coil,aspect=0.3,segment length=3.2pt,amplitude=3pt}] (0.5,3) -- (0.5,1.6);
  \draw[fill=lightgray!60] (0.1,1.1) rectangle (0.9,1.6);
  \node at (0.5,1.35) {$m$};
  \fill (0.5,1.35) circle (0.6pt);
\end{tikzpicture}
\end{center}

\begin{parts}
  \part[2] Zeichnen Sie direkt in der Skizze oben, ausgehend vom Schwerpunkt
  der Masse (markierter Punkt), die Federkraft $\vec F$ und die Gewichtskraft
  $\vec F_G$ als Pfeile ein. Wählen Sie dafür einen selbst gewählten,
  einheitlichen Kraftmassstab und beschriften Sie beide Pfeile.

  \part[1] Formulieren Sie das Kräftegleichgewicht zwischen Federkraft und
  Gewichtskraft nun als \emph{Vektorgleichung} (nicht nur als Aussage über
  die Beträge).
  \Antwortraum{1}
  \begin{solution}
  \[
    \vec F + \vec F_G = \vec 0, \qquad \text{d.\,h. } \vec F = -\vec F_G.
  \]
  Die beiden Kraftvektoren sind gleich lang, zeigen aber in \emph{entgegengesetzte}
  Richtungen; ihre Vektorsumme ist der Nullvektor.
  \end{solution}
\end{parts}
\end{taskbox}
\end{questions}

% =============================================================================
% Theorie 1: Kraft als vektorielle Groesse
% =============================================================================
\begin{theoriebox}[Kraft als vektorielle Grösse]
Bisher wurde eine Kraft meist nur über ihren \textbf{\emph{Betrag}} (die
Zahl in Newton) beschrieben. Das reicht jedoch nicht aus, um zu wissen, wie
ein Körper auf mehrere gleichzeitig wirkende Kräfte reagiert: Zwei Kräfte
von je $\SI{4}{\newton}$, die in dieselbe Richtung ziehen, wirken ganz
anders zusammen als zwei Kräfte von je $\SI{4}{\newton}$, die in
entgegengesetzte Richtungen ziehen, wie schon die Feder-Masse-Skizze oben
zeigt.

Deshalb hat jede Kraft nicht nur einen Betrag, sondern auch eine
\textbf{\emph{Richtung}}. Eine Grösse mit Betrag \emph{und} Richtung heisst
\textbf{\emph{Vektor}} und wird durch einen Pfeil dargestellt: Die Länge des
Pfeils steht (nach einem einheitlichen Kraftmassstab) für den Betrag, die
Pfeilrichtung für die Richtung der Kraft. Das Formelsymbol trägt deshalb
einen Vektorpfeil: $\vec F$, Einheit weiterhin $\si{\newton}$.

\begin{hinweisbox}[Hinweis: Kraftmassstab]
Sobald mehrere Kraftpfeile im selben Diagramm gezeichnet werden, muss für
\emph{alle} Pfeile derselbe Kraftmassstab gelten (z.\,B. $\Massstab{1}{1}$).
Sonst lässt sich das Ergebnis nachher nicht mehr mit
dem Lineal nachmessen. Ein Pfeil, der nur \glqq ungefähr passend\grqq{}
gezeichnet wird, macht die ganze Konstruktion unbrauchbar.
\end{hinweisbox}
\end{theoriebox}

% =============================================================================
% Theorie 2: Kraftaddition zweier Kraefte, Kraefteparallelogramm
% =============================================================================
\begin{theoriebox}[Kraftaddition zweier Kräfte: das Kräfteparallelogramm]
Wirken zwei Kräfte $\vec F_1$ und $\vec F_2$ am \textbf{\emph{selben
Angriffspunkt}} eines Körpers, lässt sich ihre gemeinsame Wirkung durch eine
einzige Ersatzkraft beschreiben, die \textbf{\emph{Resultierende}}
$\vec F_{\text{res}}$. Sie wird mit dem \textbf{\emph{Kräfteparallelogramm}}
konstruiert:

\begin{enumerate}[leftmargin=1.2cm,itemsep=0.2em]
  \item Beide Kräfte massstabsgetreu vom gemeinsamen Angriffspunkt aus
  einzeichnen.
  \item Das Parallelogramm ergänzen: An die Spitze von $\vec F_1$ eine Parallele
  zu $\vec F_2$ zeichnen, und an die Spitze von $\vec F_2$ eine Parallele zu
  $\vec F_1$.
  \item Die Resultierende $\vec F_{\text{res}}$ ist die \textbf{\emph{Diagonale}}
  des Parallelogramms, die durch den gemeinsamen Angriffspunkt verläuft (nicht
  die andere Diagonale!).
\end{enumerate}

\begin{center}
\begin{tikzpicture}[>=Latex,scale=1]
  \Feinraster{0}{0}{5}{5}
  \coordinate (A) at (0,0);
  \coordinate (P1) at (3,0);
  \coordinate (P2) at (0,4);
  \coordinate (R) at (3,4);
  \draw[dashed,gray] (P1) -- (R) -- (P2);
  \draw[->,thick,titleblue] (A) -- (P1) node[midway,below]{$\vec F_1$};
  \draw[->,thick,accentorange] (A) -- (P2) node[midway,left]{$\vec F_2$};
  \draw[->,ultra thick,accentpurple] (A) -- (R) node[midway,above left]{$\vec F_{\text{res}}$};
  \draw (0.7,0) arc[start angle=0,end angle=90,radius=0.7];
  \node at (0.5,0.5) {\small $\theta$};
\end{tikzpicture}
\end{center}
{\scriptsize Beispiel bei Massstab $\Massstab{1}{1}$: $F_1=\SI{3}{\newton}$
nach rechts, $F_2=\SI{4}{\newton}$ nach oben, Zwischenwinkel $\theta=90^\circ$.
Nachmessen der Diagonale mit dem Lineal ergibt $\SI{5}{\centi\meter}$, also
$F_{\text{res}}=\SI{5}{\newton}$.}
\end{theoriebox}

% =============================================================================
% Uebung: Kraftaddition mehrerer Kraefte + Alltagsbeispiel (M3 Whiteboarding)
% =============================================================================
\begin{questions}
\begin{groupbox}[Whiteboarding: Kraftaddition mehrerer Kräfte]
\question Bei mehr als zwei Kräften am selben Angriffspunkt funktioniert
das Kräfteparallelogramm nicht mehr direkt. Stattdessen verwendet man das
\textbf{\emph{Kopf-an-Schwanz-Verfahren}}: Die Kraftpfeile werden (immer im
selben Kraftmassstab) nacheinander so aneinandergehängt, dass der Schwanz
des nächsten Pfeils an der Spitze (dem \glqq Kopf\grqq{}) des vorherigen
beginnt. Die Resultierende ist dann der Pfeil vom Schwanz des ersten Pfeils
zur Spitze des letzten. Die Reihenfolge der Pfeile spielt dabei keine Rolle:
das Ergebnis bleibt gleich.

\begin{parts}
  \part[4] Bilden Sie Dreiergruppen. Wählen Sie \emph{eines} der folgenden
  Alltagsbeispiele, bei dem jeweils \textbf{\emph{drei}} Kräfte gleichzeitig
  am selben Angriffspunkt wirken, und schätzen Sie für alle drei Kräfte
  plausible Beträge und Richtungen:
  \begin{itemize}[leftmargin=1cm]
    \item Tauziehen mit drei Mannschaften, die an einem gemeinsamen
    Ring in drei unterschiedliche Richtungen ziehen.
    \item Ein Boot wird gleichzeitig vom Motor (Vortrieb), vom seitlich
    einfallenden Wind und von der Strömung des Wassers angetrieben.
    \item Ein Schlitten wird an drei Seilen von drei Personen gleichzeitig
    in leicht unterschiedliche Richtungen gezogen.
  \end{itemize}
  Konstruieren Sie auf dem Whiteboard mit dem Kopf-an-Schwanz-Verfahren die
  Resultierende aller drei Kräfte und halten Sie fest: In welche Richtung
  würde sich der Körper (Ring, Boot, Schlitten) tatsächlich in Bewegung
  setzen?

  \part[2] Übertragen Sie Ihre Konstruktion massstabsgetreu ins Feinraster
  unten (selbst gewählter Kraftmassstab, z.\,B. $\Massstab{1}{1}$)
  und beschriften Sie jeden Kraftpfeil sowie die Resultierende.
  \begin{center}
  \begin{tikzpicture}[>=Latex,scale=1.3]
    \Feinraster{0}{0}{8}{6}
  \end{tikzpicture}
  \end{center}
\end{parts}
\end{groupbox}
\end{questions}

% =============================================================================
% Kernaufgabe: Resultierende vs. Winkel (M9 Produktives Ueben, LZ3)
% =============================================================================
\begin{questions}
\begin{taskbox}[Kernaufgabe: Wie hängt die Resultierende vom Winkel ab?]
\question Zwei gleich grosse Kräfte $F_1=F_2=\SI{4}{\newton}$ greifen am
selben Punkt an. In den vier Feinrastern unten ist $F_1$ bereits waagrecht
eingezeichnet, $F_2$ jeweils im angegebenen Zwischenwinkel $\theta$ dazu
(Massstab: $\Massstab{1}{1}$).

\begin{parts}
  \part[6] Ergänzen Sie in jedem der vier Raster das Kräfteparallelogramm
  (bzw. bei $\theta=0^\circ$ und $\theta=180^\circ$ die entartete, geradlinige
  Konstruktion) und zeichnen Sie die Resultierende ein. Messen Sie deren
  Länge mit dem Lineal nach und tragen Sie den Betrag $F_{\text{res}}$ in
  die Tabelle darunter ein.

\begin{center}
\begin{tikzpicture}[>=Latex,scale=1]
  \Feinraster{0}{-1}{9}{2}
  \coordinate (A) at (0,0);
  \draw[->,thick,titleblue] (A) -- (4,0.15) node[midway,above]{\scriptsize $F_1$};
  \draw[->,thick,accentorange] (A) -- (4,-0.15) node[midway,below]{\scriptsize $F_2$};
  \node[below] at (2,-1.3) {\scriptsize $\theta=0^\circ$};
\end{tikzpicture}

\begin{tikzpicture}[>=Latex,scale=1]
  \Feinraster{0}{0}{7}{5}
  \coordinate (A) at (0,0);
  \coordinate (P1a) at (4,0);
  \coordinate (P2a) at (2,3.4641);
  \draw[->,thick,titleblue] (A) -- (P1a) node[midway,below]{\scriptsize $F_1$};
  \draw[->,thick,accentorange] (A) -- (P2a) node[midway,above left]{\scriptsize $F_2$};
  \draw (0.6,0) arc[start angle=0,end angle=60,radius=0.6];
  \node[below] at (3.5,-0.6) {\scriptsize $\theta=60^\circ$};
\end{tikzpicture}

\begin{tikzpicture}[>=Latex,scale=1]
  \Feinraster{0}{0}{5}{5}
  \coordinate (A) at (0,0);
  \draw[->,thick,titleblue] (A) -- (4,0) node[midway,below]{\scriptsize $F_1$};
  \draw[->,thick,accentorange] (A) -- (0,4) node[midway,left]{\scriptsize $F_2$};
  \node[below] at (2,-0.6) {\scriptsize $\theta=90^\circ$};
\end{tikzpicture}

\begin{tikzpicture}[>=Latex,scale=1]
  \Feinraster{-5}{-1}{5}{1}
  \coordinate (A) at (0,0);
  \draw[->,thick,titleblue] (A) -- (4,0) node[midway,above]{\scriptsize $F_1$};
  \draw[->,thick,accentorange] (A) -- (-4,0) node[midway,below]{\scriptsize $F_2$};
  \node[below] at (0,-1.3) {\scriptsize $\theta=180^\circ$};
\end{tikzpicture}
\end{center}
{\scriptsize Bei $\theta=0^\circ$ sind $F_1$ und $F_2$ zur besseren
Sichtbarkeit leicht versetzt gezeichnet. Tatsächlich zeigen beide Kräfte
exakt in dieselbe Richtung. Alle vier Raster gelten für denselben Massstab
$\Massstab{1}{1}$; nachmessen mit dem Lineal ergibt also unmittelbar den
Kraftbetrag in Newton.}

\begin{solution}
\begin{center}
\begin{tikzpicture}[>=Latex,scale=1]
  \Feinraster{0}{-1}{9}{2}
  \coordinate (A) at (0,0);
  \draw[->,thick,titleblue] (A) -- (4,0.15) node[midway,above]{\scriptsize $F_1$};
  \draw[->,thick,accentorange] (A) -- (4,-0.15) node[midway,below]{\scriptsize $F_2$};
  \draw[->,ultra thick,accentpurple] (A) -- (8,0) node[midway,above,yshift=6pt]{\scriptsize $F_{\text{res}}=\SI{8}{\newton}$};
  \node[below] at (4,-1.3) {\scriptsize $\theta=0^\circ$};
\end{tikzpicture}

\begin{tikzpicture}[>=Latex,scale=1]
  \Feinraster{0}{0}{7}{5}
  \coordinate (A) at (0,0);
  \coordinate (P1) at (4,0);
  \coordinate (P2) at ({4*cos(60)},{4*sin(60)});
  \coordinate (R) at ({4+4*cos(60)},{4*sin(60)});
  \draw[dashed,gray] (P1) -- (R) -- (P2);
  \draw[->,thick,titleblue] (A) -- (P1) node[midway,below]{\scriptsize $F_1$};
  \draw[->,thick,accentorange] (A) -- (P2) node[midway,above left]{\scriptsize $F_2$};
  \draw[->,ultra thick,accentpurple] (A) -- (R) node[midway,above]{\scriptsize $F_{\text{res}}\approx\SI{6.9}{\newton}$};
  \node[below] at (3,-0.6) {\scriptsize $\theta=60^\circ$};
\end{tikzpicture}

\begin{tikzpicture}[>=Latex,scale=1]
  \Feinraster{0}{0}{5}{5}
  \coordinate (A) at (0,0);
  \coordinate (P1) at (4,0);
  \coordinate (P2) at (0,4);
  \coordinate (R) at (4,4);
  \draw[dashed,gray] (P1) -- (R) -- (P2);
  \draw[->,thick,titleblue] (A) -- (P1) node[midway,below]{\scriptsize $F_1$};
  \draw[->,thick,accentorange] (A) -- (P2) node[midway,left]{\scriptsize $F_2$};
  \draw[->,ultra thick,accentpurple] (A) -- (R) node[midway,above left]{\scriptsize $F_{\text{res}}\approx\SI{5.7}{\newton}$};
  \node[below] at (2,-0.6) {\scriptsize $\theta=90^\circ$};
\end{tikzpicture}

\begin{tikzpicture}[>=Latex,scale=1]
  \Feinraster{-5}{-1}{5}{1}
  \coordinate (A) at (0,0);
  \draw[->,thick,titleblue] (A) -- (4,0) node[midway,above]{\scriptsize $F_1$};
  \draw[->,thick,accentorange] (A) -- (-4,0) node[midway,below]{\scriptsize $F_2$};
  \fill (0,0) circle (2pt);
  \node[below] at (0,-1.3) {\scriptsize $\theta=180^\circ$, $F_{\text{res}}=\SI{0}{\newton}$};
\end{tikzpicture}
\end{center}
Bei $\theta=0^\circ$ addieren sich die Kräfte vollständig (Vollwinkel,
Wirklinie ist dieselbe); bei $\theta=180^\circ$ heben sie sich vollständig
auf. In beiden Fällen entartet das Parallelogramm zu einer Strecke, daher
genügt das direkte Aneinanderhängen der Pfeile.
\end{solution}

  \part[4] Ergänzen Sie die Tabelle mit den vier gemessenen Werten und
  vergleichen Sie diese mit der arithmetischen Summe der beiden
  Kraftbeträge ($F_1+F_2=\SI{8}{\newton}$). Formulieren Sie eine allgemeine
  Aussage: Wie verändert sich $F_{\text{res}}$, wenn der Zwischenwinkel
  $\theta$ von $0^\circ$ auf $180^\circ$ zunimmt? Gilt
  $F_{\text{res}}=F_1+F_2$ auch bei $\theta\neq0^\circ$?

  \begin{center}
  \begin{tabular}{lcccc}
  \toprule
  $\theta$ & $0^\circ$ & $60^\circ$ & $90^\circ$ & $180^\circ$ \\
  \midrule
  $F_{\text{res}}$ (gemessen) & \rule{1.6cm}{0.4pt} & \rule{1.6cm}{0.4pt} & \rule{1.6cm}{0.4pt} & \rule{1.6cm}{0.4pt} \\
  \bottomrule
  \end{tabular}
  \end{center}
  \Antwortraum{2}
  \begin{solution}
  \begin{center}
  \begin{tabular}{lcccc}
  \toprule
  $\theta$ & $0^\circ$ & $60^\circ$ & $90^\circ$ & $180^\circ$ \\
  \midrule
  $F_{\text{res}}$ (gemessen) & $\SI{8}{\newton}$ & $\approx\SI{6.9}{\newton}$ & $\approx\SI{5.7}{\newton}$ & $\SI{0}{\newton}$ \\
  \bottomrule
  \end{tabular}
  \end{center}
  Je grösser der Zwischenwinkel $\theta$, desto \textbf{\emph{kleiner}} wird
  die Resultierende: Sie ist bei $\theta=0^\circ$ maximal (gleich der
  arithmetischen Summe $F_1+F_2$) und bei $\theta=180^\circ$ minimal (hier
  sogar null, weil $F_1=F_2$). $F_{\text{res}}=F_1+F_2$ gilt also
  \textbf{\emph{nur}} für $\theta=0^\circ$, wenn beide Kräfte in dieselbe
  Richtung zeigen. Für jeden anderen Zwischenwinkel liegt $F_{\text{res}}$
  \emph{unterhalb} dieser arithmetischen Summe: Kräfte addieren sich also
  \emph{nicht} wie gewöhnliche Zahlen, sondern wie Vektoren.
  \end{solution}
\end{parts}
\end{taskbox}
\end{questions}

\begin{formelbox}[Merksatz: Kraftaddition]
Zwei Kräfte am selben Angriffspunkt werden mit dem
\textbf{\emph{Kräfteparallelogramm}} addiert; die Resultierende ist die
Diagonale durch den gemeinsamen Angriffspunkt. Nur wenn beide Kräfte in
\emph{dieselbe} Richtung zeigen ($\theta=0^\circ$), ist die Resultierende
gleich der arithmetischen Summe der Beträge. Sonst ist sie kleiner.
\end{formelbox}

% =============================================================================
% Theorie 3: Kraftzerlegung mit dem Kraefteparallelogramm (Ampelbeispiel)
% =============================================================================
\begin{theoriebox}[Kraftzerlegung: die Umkehrung der Kraftaddition]
Manchmal ist der umgekehrte Fall gefragt: Eine Kraft ist gegeben, und man
möchte wissen, welche zwei Teilkräfte in zwei \emph{vorgegebenen} Richtungen
zusammen dieselbe Wirkung hätten. Das ist die \textbf{\emph{Kraftzerlegung}},
dieselbe Konstruktion wie bei der Kraftaddition, nur in umgekehrter
Reihenfolge:

\begin{enumerate}[leftmargin=1.2cm,itemsep=0.2em]
  \item Die gegebene Kraft $\vec F$ sowie die beiden vorgegebenen Richtungen
  vom selben Angriffspunkt aus einzeichnen.
  \item Von der Spitze von $\vec F$ aus Parallelen zu beiden Richtungen
  ziehen, bis sie diese schneiden.
  \item Die beiden Teilkräfte $\vec F_1$ und $\vec F_2$ sind die Strecken vom
  Angriffspunkt bis zu den beiden Schnittpunkten.
\end{enumerate}

\begin{hinweisbox}[Hinweis: Addition oder Zerlegung?]
Bei der \textbf{\emph{Addition}} sind die Beträge und Richtungen \emph{aller}
Kräfte bereits bekannt, gesucht ist die Resultierende. Bei der
\textbf{\emph{Zerlegung}} ist es umgekehrt: Nur \emph{eine} Kraft und die
beiden \emph{Richtungen} der gesuchten Teilkräfte sind bekannt, ihre Beträge
sind gesucht.
\end{hinweisbox}

\textbf{Beispiel:} Eine Verkehrsampel hängt an zwei schräg gespannten
Seilen. Die Gewichtskraft $\vec F_G$ der Ampel zeigt senkrecht nach unten.
Damit die Ampel in Ruhe bleibt, muss nach dem Kräftegleichgewicht
(vgl.\ Feder-Masse-Skizze, LZ1) die Vektorsumme der beiden Seilkräfte genau
$-\vec F_G$ ergeben. Deshalb wird nicht $\vec F_G$ selbst, sondern der ihr
entgegengesetzte Vektor $-\vec F_G$ in die beiden vorgegebenen
Seilrichtungen zerlegt: Die so gefundenen Teilkräfte $\vec F_1$, $\vec F_2$
sind die Seilkräfte, mit denen die Seile tatsächlich an der Ampel ziehen und
sie so im Gleichgewicht halten.

\begin{center}
\begin{tikzpicture}[>=Latex,scale=1.6]
  \coordinate (A) at (0,0);
  \coordinate (S1) at (-3,2.2);
  \coordinate (S2) at (3,1.4);
  \coordinate (P1) at (-2.5,1.8333);
  \coordinate (P2) at (2.5,1.1667);
  \coordinate (M) at (0,3);
  \draw[very thick,gray] (A) -- (S1);
  \draw[very thick,gray] (A) -- (S2);
  \draw[->,ultra thick,titleblue] (A) -- (0,-3) node[midway,right]{$\vec F_G$};
  \draw[->,dashed,darkgray] (A) -- (M) node[midway,left]{\scriptsize $-\vec F_G$};
  \draw[dashed,gray] (M) -- (P1);
  \draw[dashed,gray] (M) -- (P2);
  \draw[->,thick,accentorange] (A) -- (P1) node[midway,above]{$\vec F_1$};
  \draw[->,thick,accentpurple] (A) -- (P2) node[midway,above right]{$\vec F_2$};
\end{tikzpicture}
\end{center}
{\scriptsize Zerlegung von $-\vec F_G$ (gestrichelt) in zwei Teilkräfte
$\vec F_1$, $\vec F_2$ entlang der vorgegebenen Seilrichtungen (nach
PhysikLibre, Bild 4.17). Die grauen Linien markieren die Seilrichtungen, die
kurzen gestrichelten Linien die Parallelen, die $\vec F_1$ und $\vec F_2$
festlegen.}
\end{theoriebox}

% =============================================================================
% Uebung: Eigene Zerlegung + freies Alltagsbeispiel (Partnerarbeit, LZ4/LZ5)
% =============================================================================
\begin{questions}
\begin{groupbox}[Partnerarbeit: Kraftzerlegung selbst konstruieren]
\question
\begin{parts}
  \part[4] Ein Wäscheseil ist zwischen zwei Balkonen gespannt. In der Mitte
  hängt ein nasses Wäschestück, das eine Gewichtskraft von $F_G=\SI{6}{\newton}$
  senkrecht nach unten ausübt. Das Seil verläuft auf beiden Seiten des
  Wäschestücks schräg nach oben, wie im Feinraster unten angedeutet.
  Zerlegen Sie $\vec F_G$ grafisch in die beiden Teilkräfte $\vec F_1$,
  $\vec F_2$ entlang der beiden Seilrichtungen und messen Sie deren Beträge.

  \begin{center}
  \begin{tikzpicture}[>=Latex,scale=1]
    \Feinraster{-4}{-7}{4}{3.5}
    \coordinate (A) at (0,0);
    \draw[very thick,gray] (-3.5,1.4) -- (A) -- (3.5,2.6);
    \draw[->,ultra thick,titleblue] (A) -- (0,-6) node[midway,right]{$\vec F_G=\SI{6}{\newton}$};
  \end{tikzpicture}
  \end{center}
  \Antwortraum{2}
  \begin{solution}
  Konstruktion analog zum Ampelbeispiel: Von der Spitze von $\vec F_G$ aus
  Parallelen zu den beiden Seilrichtungen ziehen; die Abschnitte vom
  Angriffspunkt bis zu den Schnittpunkten sind $\vec F_1$ und $\vec F_2$. Da
  die beiden Seilrichtungen hier unterschiedlich steil sind, sind auch
  $F_1$ und $F_2$ unterschiedlich gross: das steilere Seilstück trägt den
  grösseren Anteil von $F_G$.
  \end{solution}

  \part[4] Wählen Sie zu zweit \emph{ein} eigenes Alltagsbeispiel, entweder
  für Kraftaddition oder für Kraftzerlegung, nicht notwendigerweise eines der
  bereits behandelten. Skizzieren Sie die Situation kurz, zeichnen Sie die
  beteiligten Kräfte massstabsgetreu in einem eigenen Feinraster ein und
  konstruieren Sie Resultierende bzw. Teilkräfte.
  \Antwortraum{5}
\end{parts}
\end{groupbox}
\end{questions}

\end{document}
